% ---------------------------------------------------------------------
\subsubsection{Requisiti di Sistema}
\label{subsec:requisiti-sistema}
% ---------------------------------------------------------------------

Questa sezione dettaglia i requisiti specifici che il sistema \productname{} deve soddisfare. 
I requisiti sono definiti da una prospettiva esterna, descrivendo le interazioni del sistema con i suoi utenti e con altri sistemi.

% ---------------------------------------------------------------------
\paragraph{Requisiti delle Interfacce Esterne}
\label{subsubsec:interfacce-esterne}
% ---------------------------------------------------------------------

\subparagraph{Interfaccia Utente (User Interface)}
Il sistema dovrà fornire un'interfaccia utente grafica (GUI) web-based, accessibile tramite i browser moderni. 
L'interfaccia dovrà essere responsive, garantendo l'usabilità su schermi di diverse dimensioni, dai dispositivi mobili
(minimo 320px di larghezza) ai desktop. 
Tutte le funzionalità esposte agli utenti dovranno essere accessibili tramite tale interfaccia.

\subparagraph{Interfaccia Software (Software Interface)}
Il sistema dovrà esporre un'interfaccia di programmazione (API) basata sullo stile architetturale REST 
per consentire la comunicazione tra il componente front-end e il back-end. 
Queste API dovranno utilizzare il formato JSON per lo scambio di dati e seguire 
le convenzioni standard del protocollo HTTP per le operazioni (GET, POST, PUT, DELETE). 
Questo garantirà l'interoperabilità e la possibilità di sviluppare client alternativi in futuro.

\subparagraph{Interfaccia di Comunicazione (Communications Interface)}
Tutta la comunicazione tra il client e il server che trasporta dati sensibili 
(incluse le credenziali di autenticazione e i dati delle issue) 
dovrà avvenire tramite protocolli di comunicazione sicuri (HTTPS).

% ---------------------------------------------------------------------
\paragraph{Requisiti Prestazionali}
\label{subsubsec:requisiti-prestazionali}
% ---------------------------------------------------------------------
Il sistema dovrà rispettare i requisiti di performance specificati nella sezione dei 
Requisiti Non Funzionali (REQ-NFR-01), garantendo tempi di risposta rapidi per le operazioni interattive.

% ---------------------------------------------------------------------
\paragraph{Requisiti del Database Logico}
\label{subsubsec:database-logico}
% ---------------------------------------------------------------------
Il sistema dovrà garantire la persistenza, la consistenza e l'integrità dei dati. 
Dovrà gestire le seguenti entità logiche principali e le loro relazioni: 
Utenti, Progetti, Issue, Commenti, Etichette, Allegati e le associazioni tra di esse 
(es. Utenti membri di Progetti, Issue appartenenti a Progetti). 
Il sistema dovrà garantire che le operazioni sul database siano transazionali per prevenire la corruzione dei dati in caso di fallimenti.
\newpage